% !TEX root = ../main.tex
\chapter{Shape, View, and Broadcasting Semantics}
\label{ch:shape_view}

Shape manipulation operations change a tensor's logical layout
without (in most cases) touching its element values. They are the
operations that turn a $(B, C, H, W)$ activation into a
$(B \cdot C, H \cdot W)$ matrix for a fully-connected layer, that
transpose a weight matrix for a backward pass, that broadcast a
bias vector across a batch. The chapter describes the operations
themselves, the view-vs-copy semantics that determine when memory
is touched, and the broadcasting rules that the binary kernels
inherit.

\chref{ch:tensor_model} already introduced the storage model that
makes views cheap: shape, strides, and storage offset are
per-\code{TensorImpl} fields, the underlying \code{Storage} is
shared via reference count, and a view operation produces a new
\code{TensorImpl} with adjusted metadata. This chapter assumes that
background and focuses on the user-facing surface.

\secref{sec:shape_view:reshape} treats reshape and view.
\secref{sec:shape_view:permute} treats transpose and permute.
\secref{sec:shape_view:squeeze} treats squeeze and unsqueeze.
\secref{sec:shape_view:expand} treats expand, broadcast\_to, and
repeat. \secref{sec:shape_view:concat} treats concatenation and
stacking. \secref{sec:shape_view:split} treats split and chunk.
\secref{sec:shape_view:broadcast} treats the broadcasting rules.
\secref{sec:shape_view:contiguity} closes with the contiguity
invariants.

\tabref{tab:shape_view:ops_summary} classifies every shape op by
whether it returns a view (free) or a fresh tensor (allocates),
and by what its backward does. The remainder of the chapter
expands on each.

\begin{table}[t]
\centering
\footnotesize
\setlength{\tabcolsep}{4pt}
\begin{adjustbox}{max width=\textwidth, center}
\begin{tabular}{lllll}
\toprule
Op & Always view? & Output rank vs input & Backward & Notes \\
\midrule
\code{reshape}    & if contiguous     & equal/diff & inverse reshape  & may copy if non-contig \\
\code{view}       & yes (errors else) & equal/diff & inverse reshape  & strict version \\
\code{transpose}  & yes               & equal      & inverse transpose & swap two dims \\
\code{permute}    & yes               & equal      & inverse permute   & general permutation \\
\code{squeeze}    & yes               & smaller    & \code{unsqueeze}  & remove size-1 dims \\
\code{unsqueeze}  & yes               & larger     & \code{squeeze}    & insert size-1 dim \\
\code{slice}      & yes               & equal      & zero-pad          & possibly negative stride \\
\code{flip}       & yes               & equal      & inverse flip      & negative stride \\
\code{expand}     & yes (stride 0)    & equal      & sum over expanded & read-only \\
\code{broadcast\_to} & yes (stride 0) & possibly larger & sum over broadcast & alias for expand \\
\code{repeat}     & no (allocates)    & equal      & sum over repeats  & physical copy \\
\code{cat}        & no (allocates)    & equal      & \code{split}      & concatenate along axis \\
\code{stack}      & no (allocates)    & larger     & index along new dim & insert new axis \\
\code{split}      & yes (each slice)  & equal      & \code{cat}        & multiple views \\
\code{chunk}      & yes (each slice)  & equal      & \code{cat}        & equal-size split \\
\code{contiguous} & no (if non-contig) & equal     & identity          & materialise unit stride \\
\bottomrule
\end{tabular}
\end{adjustbox}
\caption[Shape-op summary.]{Classification of \Lucid's shape ops
by whether they return a view (cost: only metadata) or a fresh
tensor (cost: data copy). The view-vs-copy distinction is the
single most useful fact about a shape op for predicting its
performance impact.}
\label{tab:shape_view:ops_summary}
\end{table}

\section{Reshape and View}
\label{sec:shape_view:reshape}

\code{reshape(new\_shape)} returns a tensor of the requested shape
with the same number of elements. If the source is contiguous, the
operation returns a view (no copy); if it is not, the operation
materialises a contiguous copy first and then returns a view of
the copy.

The distinction matters when the user expects a view. \Lucid\
mirrors \refframework\ in providing two related operations:

\begin{itemize}
  \item \code{x.view(new\_shape)}: returns a view, or raises
    \code{LucidError} if the source is not contiguous in the
    relevant axes. Used when the caller knows a view is possible
    and wants the error otherwise.
  \item \code{x.reshape(new\_shape)}: returns a view if possible,
    a copy otherwise. The conservative default.
\end{itemize}

\subsection{The -1 placeholder}
\label{ssec:shape_view:minus_one}

\code{reshape} and \code{view} accept \code{-1} in at most one
dimension as a wildcard. The size of that dimension is inferred
from the total number of elements and the other declared sizes.

\code{x.reshape(-1, 16)} on a tensor of \(64 \cdot 32\) elements
produces a \((128, 16)\) tensor; the wildcard is inferred as
\(128\). A \code{LucidError} is raised if the inference fails (the
declared sizes do not divide the total).

\subsection{The compatible-stride check}
\label{ssec:shape_view:stride_check}

A reshape can be a view if and only if the input's strides are
compatible with the output's strides for some contiguous mapping.
The check is straightforward but not trivial: the input's strides
must factor into runs that align with the output's dimensions. For
a contiguous input, every reshape is compatible; for an arbitrary
input, some reshapes are compatible (those that preserve the
order of dimensions) and some are not (those that interleave).

The implementation in \code{ops/gfunc/Shape.cpp} walks both shape
vectors and verifies the factoring. If incompatible, the operation
materialises a contiguous copy.

\section{Transpose and Permute}
\label{sec:shape_view:permute}

\code{transpose(dim0, dim1)} swaps two axes; \code{permute(*dims)}
applies an arbitrary permutation. Both always return views: shape
and strides are permuted accordingly, the underlying storage is
unchanged.

\code{x.transpose(0, 1)} on a $(B, T, D)$ tensor returns a
$(T, B, D)$ tensor without data motion. The resulting view is
typically not contiguous in C order; a subsequent operation that
requires contiguity will trigger materialisation.

\subsection{The .T shortcut}
\label{ssec:shape_view:t_shortcut}

For rank-2 tensors, \code{x.T} is shorthand for
\code{x.transpose(0, 1)}. For higher-rank tensors, \code{x.T}
reverses all dimensions: \(x_{i,j,k} \to x_{k,j,i}\). Matching
\refframework, \Lucid\ deprecates \code{.T} on rank \(\neq 2\)
tensors with a warning; the recommended replacement is
\code{x.permute(2, 1, 0)} or \code{x.mT} for matrix-transpose
semantics on the trailing two axes.

\subsection{The mT shortcut for batched matrices}
\label{ssec:shape_view:mt_shortcut}

\code{x.mT} (matrix transpose) swaps the last two axes only,
leaving leading batch dimensions untouched. For a $(B, M, N)$
batched matrix, \code{x.mT} returns $(B, N, M)$. This is the
common case for \code{matmul} workflows and is preferable to
\code{.T} for rank > 2.

\section{Squeeze and Unsqueeze}
\label{sec:shape_view:squeeze}

\code{squeeze(dim=None)} removes dimensions of size 1. With
\code{dim} specified, only that dimension is removed (and only if
it has size 1; otherwise no-op). Without \code{dim}, every
size-1 dimension is removed.

\code{unsqueeze(dim)} inserts a new dimension of size 1 at
position \code{dim}. The result has rank one greater than the
input.

Both are views: only the shape and strides change.

\subsection{Why explicit dim for squeeze}
\label{ssec:shape_view:explicit_dim}

The unconditional \code{squeeze()} is convenient but fragile:
a tensor that happened to have a size-1 batch dimension would
lose its batch axis silently, and a downstream operation
expecting a batch axis would mis-broadcast. The recommended idiom
in production code is to pass \code{dim} explicitly:
\code{x.squeeze(dim=1)} squeezes axis 1 if it is size 1, otherwise
returns \code{x} unchanged. This is harder to misuse.

\section{Expand, Broadcast\_to, and Repeat}
\label{sec:shape_view:expand}

These three operations all enlarge a tensor by replicating along
some axis, but they differ in whether the replication is virtual
or physical.

\subsection{Expand: virtual replication}
\label{ssec:shape_view:expand_virtual}

\code{expand(new\_shape)} requires that any dimension being
enlarged has size 1 in the source. The result is a view with the
relevant strides set to 0 --- the data is not copied; reads along
the enlarged axis return the same byte.

\code{x.expand(B, C, H, W)} on a $(1, C, 1, W)$ tensor returns a
$(B, C, H, W)$ view that uses zero memory beyond the original.
This is the same mechanism the binary kernels use for broadcasting
(\chref{ch:binary_functions}).

\subsection{Broadcast\_to: convenience wrapper}
\label{ssec:shape_view:broadcast_to}

\code{broadcast\_to(new\_shape)} is a friendlier alias for
\code{expand} that accepts the NumPy-style broadcasting rule (any
dimension of size 1 may be expanded). The two are functionally
equivalent; \code{broadcast\_to} is the NumPy compatibility name.

\subsection{Repeat: physical replication}
\label{ssec:shape_view:repeat_physical}

\code{repeat(reps)} copies the tensor data \code{reps} times along
each dimension. The result is a fresh contiguous tensor with
duplicated data.

\code{x.repeat(2, 3, 1, 1)} on a $(1, 1, H, W)$ tensor returns a
$(2, 3, H, W)$ tensor that occupies six times the original
storage. \code{repeat} is what the user wants when downstream code
needs an actual replicated buffer (e.g., for an in-place operation
that would otherwise alias).

\subsection{When to use which}
\label{ssec:shape_view:expand_vs_repeat}

The choice is between memory and write-safety:

\begin{itemize}
  \item \code{expand} is free in memory and time, but the result
    aliases the source. Writing to the expanded tensor writes the
    same byte multiple times --- usually a bug.
  \item \code{repeat} allocates and copies, but the result is a
    fresh tensor that can be written to safely.
\end{itemize}

The rule of thumb: \code{expand} for read-only use cases
(broadcasting in a binary op), \code{repeat} for cases that need
to write or that pass the result to a function that may write.

\section{Concatenation and Stacking}
\label{sec:shape_view:concat}

\code{cat(tensors, dim)} concatenates a sequence of tensors along
\code{dim}, requiring matching sizes on all other dimensions.
\code{stack(tensors, dim)} creates a new dimension at \code{dim}
and stacks along it, requiring matching shapes on all dimensions.

Both operations allocate: the result is a fresh contiguous tensor
holding copies of the inputs. There is no view-version, because
the input tensors typically have non-contiguous storage that
cannot be laid out as a single contiguous block without copying.

\subsection{Backward through cat and stack}
\label{ssec:shape_view:concat_back}

The backward through \code{cat} splits the gradient along
\code{dim} and routes each piece to the corresponding input. The
backward through \code{stack} indexes the gradient along the new
dimension. Both are exact inverses of the forward.

\section{Split and Chunk}
\label{sec:shape_view:split}

\code{split(x, sizes, dim)} splits \code{x} into chunks of the
given sizes along \code{dim}. \code{chunk(x, n\_chunks, dim)}
splits into \(n\) approximately equal chunks.

Both are view operations on a contiguous input: each returned
chunk is a view into a section of the source. The backward is the
inverse of \code{cat}.

\section{Broadcasting Rules}
\label{sec:shape_view:broadcast}

Broadcasting governs how two tensors of different shapes are
combined in elementwise operations
(\chref{ch:binary_functions}). The rules:

\begin{enumerate}
  \item Pad the shorter shape with leading 1s.
  \item For each dimension, the output size is the maximum of the
    two inputs.
  \item The broadcast is legal if and only if, at each dimension,
    the two input sizes are either equal or one of them is 1.
\end{enumerate}

A \code{LucidError} with a structured shape comparison is raised
on incompatible shapes.

\subsection{NumPy compatibility}
\label{ssec:shape_view:numpy_compat}

The broadcasting rules are bit-for-bit compatible with NumPy and
\refframework. Tensors of identical shape broadcast trivially;
tensors with one being a scalar broadcast trivially; tensors with
matching trailing dimensions broadcast naturally. Edge cases
(0-dimensional tensors, tensors with one size-0 dimension) follow
the NumPy convention.

\subsection{The output-shape function}
\label{ssec:shape_view:output_shape}

The pure-Python utility \code{lucid.broadcast\_shapes(*shapes)}
returns the broadcast output shape for a list of input shapes
without performing any operation. It is useful for asserting
broadcast compatibility at module-construction time:

\begin{lstlisting}[style=lucid,language=LucidPython]
out_shape = lucid.broadcast_shapes(
    (B, 1, H, W),
    (1, C, 1, W),
)
# out_shape == (B, C, H, W)
\end{lstlisting}

The function raises if the shapes are incompatible.

\section{Contiguity Invariants}
\label{sec:shape_view:contiguity}

Contiguity --- whether a tensor's strides match the canonical
C-order layout --- is a derived property that the dispatcher
consults to choose between fast and slow kernel paths.

\subsection{The contiguity predicate}
\label{ssec:shape_view:contig_pred}

\code{x.is\_contiguous()} returns \code{True} if the tensor's
strides match $(d_1 d_2 \cdots d_{n-1},\ d_2 \cdots d_{n-1},\
\ldots,\ 1)$ for shape $(d_0, d_1, \ldots, d_{n-1})$. The check
is \(O(rank)\); the result is cached on the \code{TensorImpl} for
subsequent queries.

\subsection{Operations that preserve contiguity}
\label{ssec:shape_view:contig_preserve}

\code{reshape} (on a contiguous source), \code{squeeze},
\code{unsqueeze}, and trailing-axis \code{slice} preserve
contiguity. \code{transpose}, \code{permute}, leading-axis
\code{slice}, and \code{flip} break contiguity.

A tensor that has been operated on by a chain of view operations
may have non-canonical strides. The dispatcher reads
\code{is\_contiguous()} and either runs a stride-aware kernel
(common case for elementwise ops) or materialises via
\code{.contiguous()} (when the kernel requires unit stride).

\subsection{The .contiguous() materialisation}
\label{ssec:shape_view:contiguous_op}

\code{x.contiguous()} returns \code{x} unchanged if it is already
contiguous, or a new contiguous copy otherwise. The materialisation
is one read-write pass over the tensor's elements.

The operation is most often invoked internally by the dispatcher
or by the \code{SharedStorage} bridge
(\chref{ch:unified_memory}); user code rarely calls it directly.
When it does, the typical pattern is before a numpy bridge:

\begin{lstlisting}[style=lucid,language=LucidPython]
arr = x.transpose(0, 2).contiguous().numpy()
\end{lstlisting}

The explicit \code{contiguous} ensures the numpy buffer is
unit-stride before the bridge.

\section{Worked Stride Examples}
\label{sec:shape_view:stride_examples}

The shape, stride, and storage-offset triple is the single most
useful mental model for predicting whether a view operation is
free. This section works through three representative cases.

\subsection{Example: transpose of a contiguous matrix}
\label{ssec:shape_view:transpose_example}

Start with \code{x = lucid.arange(12).reshape(3, 4)}. The metadata:

\begin{itemize}
  \item shape: $(3, 4)$.
  \item strides: $(4, 1)$ --- step 4 to advance one row, step 1 to
    advance one column.
  \item storage offset: 0.
  \item storage: 12 contiguous int64 values.
\end{itemize}

Now \code{y = x.transpose(0, 1)}:

\begin{itemize}
  \item shape: $(4, 3)$ --- dimensions swapped.
  \item strides: $(1, 4)$ --- strides swapped accordingly.
  \item storage offset: 0 (unchanged).
  \item storage: same 12 values.
\end{itemize}

Reading $y_{i,j}$ uses offset $i \cdot 1 + j \cdot 4 = i + 4j$, which
indexes the underlying contiguous storage at the original
$x_{j, i}$ position. The data is unchanged; only the iteration
order is.

The result is non-contiguous: $y$'s strides $(1, 4)$ do not match
canonical $(3, 1)$ for shape $(4, 3)$. Any subsequent
\code{reshape(12)} on $y$ will materialise.

\subsection{Example: slice with non-trivial step}
\label{ssec:shape_view:slice_example}

\code{z = x[:, ::2]}: every other column.

\begin{itemize}
  \item shape: $(3, 2)$ --- 4 columns become 2.
  \item strides: $(4, 2)$ --- column stride doubled.
  \item storage offset: 0.
  \item storage: same 12 values, only every other accessed.
\end{itemize}

Reading $z_{i,j}$ uses offset $i \cdot 4 + j \cdot 2$, which
yields the original $x_{i, 2j}$ positions. Half the storage is
unused; the data is unchanged.

\subsection{Example: expand to broadcast}
\label{ssec:shape_view:expand_example}

\code{b = lucid.tensor([1.0, 2.0, 3.0]).expand(4, 3)}: replicate
the 3-vector 4 times.

\begin{itemize}
  \item shape: $(4, 3)$.
  \item strides: $(0, 1)$ --- stride 0 along the replicated axis.
  \item storage offset: 0.
  \item storage: only 3 values.
\end{itemize}

Reading $b_{i,j}$ uses offset $i \cdot 0 + j \cdot 1 = j$, which
returns the same 3 values regardless of $i$. The replication is
virtual; the storage is unchanged.

Writing to $b$ is dangerous: every $b_{i, j}$ for the same $j$ aliases
the same byte. A scatter-style write would produce
race-conditioned results. \Lucid\ does not prevent the write
(the autograd version check catches the most common case), but the
convention is to \code{.clone()} before any write to an expanded
tensor.

\section{Contiguity Variants in Practice}
\label{sec:shape_view:contig_variants}

The contiguity predicate is more nuanced than a single boolean.
\tabref{tab:shape_view:contig_kinds} enumerates the variants the
dispatcher and kernels distinguish.

\begin{table}[t]
\centering
\small
\begin{tabular}{lll}
\toprule
Variant & Definition & Lucid use \\
\midrule
fully contiguous   & strides match canonical C-order & fast path everywhere \\
inner contiguous   & innermost axis has unit stride & sufficient for elementwise \\
outer contiguous   & strides decrease monotonically & sufficient for reductions on outer \\
zero stride        & one or more axes have stride 0 & broadcast-expanded; cannot write \\
negative stride    & one or more strides are negative & from \code{flip}; usually materialise \\
non-contiguous     & none of the above & materialise before kernel \\
\bottomrule
\end{tabular}
\caption[Contiguity variants.]{The dispatcher and kernel templates
distinguish several flavours of (non-)contiguity. The decision of
whether to materialise via \code{.contiguous()} or to dispatch a
strided kernel depends on which variant applies and which kernel
is being called.}
\label{tab:shape_view:contig_kinds}
\end{table}

\section{Broadcasting Pitfalls}
\label{sec:shape_view:pitfalls}

A short catalogue of broadcasting bugs we have seen in user code.

\subsection{Pitfall: silent dimension insertion}
\label{ssec:shape_view:silent_insert}

The right-alignment broadcasting rule means that combining a
$(B, T, D)$ tensor with a $(D,)$ tensor works, but combining
$(B, D, T)$ with $(D,)$ also works --- broadcasting the $D$ tensor
along $B$ and the inner $T$ axis. The latter is almost certainly a
bug (the dimensions don't line up semantically), but the broadcast
is mathematically legal.

The mitigation: when in doubt, use \code{None}-indexing to make
the alignment explicit. \code{x + bias[None, :, None]} is
unambiguous; \code{x + bias} requires the reader to think.

\subsection{Pitfall: writing through an expanded view}
\label{ssec:shape_view:expand_write}

\code{x.expand(4, 3)[2, 1] = 0.5} writes to the same storage byte
that every other index along the first axis reads. The result is
that all four ``copies'' of position $(\cdot, 1)$ become $0.5$, not
just $(2, 1)$.

The mitigation: \code{x.expand(4, 3).clone()[2, 1] = 0.5}.

\subsection{Pitfall: cat versus stack confusion}
\label{ssec:shape_view:cat_vs_stack}

\code{cat([a, b], dim=0)} requires $a$ and $b$ to have matching
shapes except along dim 0. \code{stack([a, b], dim=0)} requires
matching shapes on every dim and inserts a new dim 0 of size 2.
The two are easy to confuse and produce different output ranks.

The mitigation: choose by intent --- ``concatenate along an existing
axis'' is \code{cat}; ``create a new axis'' is \code{stack}.

\section{Summary and Forward References}
\label{sec:shape_view:summary}

Shape manipulation operations in \Lucid\ fall into three groups:
view operations (\code{reshape}, \code{view}, \code{transpose},
\code{permute}, \code{squeeze}, \code{unsqueeze}, \code{slice},
\code{flip}, \code{expand}) that adjust metadata without copying
data; copy operations (\code{cat}, \code{stack}, \code{repeat})
that allocate; and split operations (\code{split}, \code{chunk})
that return views into a contiguous source. The broadcasting rules
underlying the binary kernels are NumPy-compatible. The contiguity
predicate is the dispatcher's input for fast-vs-slow path
selection.

The next chapter, \chref{ch:indexing}, treats indexing and
gather/scatter --- the operations that select or update
\emph{values} from a tensor based on index tensors.
\chref{ch:composite_ops} closes Part III with the pure-Python
composite layer.

\begin{lucidnote}
For the user: most shape operations are free. The expression
\code{x.transpose(1, 2).reshape(B, -1)} produces a transposed
matrix-like view; the transpose is free (stride-only), the reshape
copies if the transposed view is non-contiguous. The decision to
\code{.contiguous()} explicitly is the user's optimisation
opportunity in pipelines that read the result repeatedly.
\end{lucidnote}
